\subsection{Zero-modes in magnetized $T^6/\mathbb{Z}_N$ orbifold models through $Sp(6,\mathbb{Z})$ modular symmetry --- arXiv:2305.16709}

\textbf{Authors:} Shota Kikuchi, Tatsuo Kobayashi, Kaito Nasu, Shohei Takada, Hikaru Uchida

\paragraph{主な主張}
磁化した非因子化$T^6/\mathbb{Z}_N$のフェルミオン零モードを解析し、調べた$\mathbb{Z}_{12}$模型では四次元$\mathcal{N}=1$超対称性を保つ三世代を実現できる一方、$\mathbb{Z}_7$模型では実現できないと示す~\cite{Kikuchi:2023awm}。

\paragraph{新規性}
$Sp(6,\mathbb{Z})$のモジュラー変換を用い、非因子化六次元オービフォールドにおける零モードの縮退度と世代数を求める。

\paragraph{位置付け}
磁化コンパクト化の世代構造を、高次元トーラスの幾何とモジュラー対称性から調べる研究。

\paragraph{コメント（読書メモ）}
本人の読書メモ：$T^2$で用いる$SL(2,\mathbb{Z})$に対し、本論文の高次元トーラスの枠組みでは$Sp(2g,\mathbb{Z})$を用い、$T^6$は$g=3$に対応する。$T^4$や$T^6$へ拡張すると、複素構造を記述するモジュライも複数現れる点に注目した。
